\begin{thebibliography}{99}
\phantomsection
\addcontentsline{toc}{section}{Références}
\bibitem{bennett2003} C. H. Bennett, ``Notes on Landauer's principle,
reversible computation, and Maxwell's Demon'', \emph{Studies in History
and Philosophy of Modern Physics} \textbf{34}, 501--510 (2003).
\href{https://doi.org/10.1016/S1355-2198(03)00039-X}{doi:10.1016/S1355-2198(03)00039-X}.
\bibitem{feynman1948} R. P. Feynman, ``Space-Time Approach to Non-Relativistic Quantum Mechanics'', \emph{Reviews of Modern Physics} \textbf{20}, 367--387 (1948). \href{https://doi.org/10.1103/RevModPhys.20.367}{doi:10.1103/RevModPhys.20.367}.
\bibitem{mustovicary} B. Musto et J. Vicary, ``Quantum Latin Squares and Unitary Error Bases'', \emph{Quantum Information and Computation} \textbf{16}, 1318--1332 (2016). \href{https://doi.org/10.26421/QIC16.15-16-4}{doi:10.26421/QIC16.15-16-4}.
\bibitem{delascuevas2020} G. De las Cuevas, T. Drescher et T. Netzer, ``Quantum Magic Squares: Dilations and Their Limitations'', \emph{Journal of Mathematical Physics} \textbf{61}, 111704 (2020). \href{https://doi.org/10.1063/5.0022344}{doi:10.1063/5.0022344}.
\bibitem{delascuevas2023} G. De las Cuevas, T. Netzer et I. Valentiner-Branth, ``Magic Squares: Latin, Semiclassical, and Quantum'', \emph{Journal of Mathematical Physics} \textbf{64}, 022201 (2023). \href{https://doi.org/10.1063/5.0127393}{doi:10.1063/5.0127393}.
\bibitem{codata2022} P. J. Mohr, D. B. Newell, B. N. Taylor et E. Tiesinga, ``CODATA Recommended Values of the Fundamental Physical Constants: 2022'', \emph{Reviews of Modern Physics} \textbf{97}, 025002 (2025). \href{https://doi.org/10.1103/RevModPhys.97.025002}{doi:10.1103/RevModPhys.97.025002}. \href{https://physics.nist.gov/cuu/Constants/Table/allascii.txt}{Tableau des constantes}.
\bibitem{flm} I. B. Frenkel, J. Lepowsky et A. Meurman, \emph{Vertex Operator Algebras and the Monster}, Pure and Applied Mathematics, vol. 134, Academic Press, 1988.
\bibitem{borcherds} R. E. Borcherds, ``Monstrous Moonshine and Monstrous Lie Superalgebras'', \emph{Inventiones Mathematicae} \textbf{109}, 405--444 (1992). \href{https://doi.org/10.1007/BF01232032}{doi:10.1007/BF01232032}.
\bibitem{conwaysloane} J. H. Conway et N. J. A. Sloane, \emph{Sphere Packings, Lattices and Groups}, 3\textsuperscript{e} éd., Grundlehren der mathematischen Wissenschaften 290, Springer, 1999. \href{https://doi.org/10.1007/978-1-4757-6568-7}{doi:10.1007/978-1-4757-6568-7}.
\bibitem{bipm2018} Conférence générale des poids et mesures, \emph{Resolution 1 of the 26th CGPM: On the revision of the International System of Units}, 2018. \href{https://www.bipm.org/en/committees/cg/cgpm/26-2018/resolution-1}{Résolution officielle}. \href{https://doi.org/10.59161/CGPM2018RES1E}{doi:10.59161/CGPM2018RES1E}.
\bibitem{duffokunveneziano} M. J. Duff, L. B. Okun et G. Veneziano, ``Trialogue on the number of fundamental constants'', \emph{JHEP} 03 (2002), 023. \href{https://arxiv.org/abs/physics/0110060}{arXiv:physics/0110060}.
\bibitem{holst1996} S. Holst, ``Barbero's Hamiltonian derived from a generalized Hilbert--Palatini action'', \emph{Physical Review D} \textbf{53} (1996), 5966--5969. \href{https://arxiv.org/abs/gr-qc/9511026}{arXiv:gr-qc/9511026}.
\bibitem{bipmbrochure} Bureau international des poids et mesures, \emph{The International System of Units}, 9\textsuperscript{e} édition et mises à jour. \href{https://www.bipm.org/en/publications/si-brochure}{Édition officielle}.
\bibitem{hmtg26radial} \emph{Rigidité simultanée du lecteur radial}.
Développement HMT--MD du 26 septembre 2026, reproduit et démontré
dans la section~\ref{g26:sec:rigidez-radial-es} ; source et vérificateurs
conservés dans le supplément de reproduction de cette édition.
\bibitem{hmtg26metric} \emph{Composition métrique, mémoire et couplage}.
Développement HMT--MD du 25 septembre 2026 ; démonstration du transport
avec métrique globale et élimination exacte de la mémoire, incorporée dans
\ref{g26:sec:rigidez-radial-es}.
\bibitem{hmtg26counter} \emph{Substitution du contre-angle dans le couplage
gravitationnel radial}. Développement HMT--MD du 25 septembre 2026 ;
équations et vérificateur de substitution conservés dans le supplément
\texttt{reproduccion\_gravitatoria\_20260926}.
\end{thebibliography}
